\documentclass[11pt,a4paper]{article}

\usepackage[T1]{fontenc}
\usepackage[utf8]{inputenc}
\usepackage[french]{babel}
\usepackage{lmodern}
\usepackage{geometry}
\usepackage{hyperref}
\usepackage{enumitem}
\usepackage{array}
\usepackage{longtable}
\usepackage{xcolor}
\usepackage{fancyvrb}

\geometry{margin=2.2cm}
\hypersetup{colorlinks=true,linkcolor=blue,urlcolor=blue}
\setlist[itemize]{noitemsep,topsep=4pt}
\setlist[enumerate]{noitemsep,topsep=4pt}

\title{\textbf{Rapport de Projet -- Développement Web ING1}\\
\large Plateforme Maison Intelligente, performante et responsive}
\author{Groupe projet}
\date{Année universitaire 2025--2026}

\begin{document}

\maketitle

\section*{Membres du groupe}
\begin{itemize}
  \item Khayem Ben Ghorbel
  \item Rindra (nom complet à compléter)
  \item Ziyad (nom complet à compléter)
  \item Anastasia (nom complet à compléter)
  \item Valerie (nom complet à compléter)
\end{itemize}

\tableofcontents
\newpage

\section{Introduction et contexte}

Le projet s'inscrit dans le cadre du module ING1 \textit{Projet Développement Web 2025--2026}.
Le sujet consiste à réaliser une plateforme intelligente autour de l'Internet des Objets,
avec un thème libre et des modules fonctionnels imposés.

Nous avons choisi le thème \textbf{Maison Intelligente}. L'objectif est de centraliser :
\begin{itemize}
  \item la consultation d'informations (objets, pièces, services),
  \item la visualisation pour les membres,
  \item la gestion opérationnelle avancée,
  \item l'administration de la plateforme.
\end{itemize}

Le projet répond aux contraintes principales du cahier des charges :
\begin{itemize}
  \item utilisation de frameworks,
  \item base de données réelle,
  \item moteur de recherche avec au moins deux filtres,
  \item dépôt Git avec commits réguliers,
  \item design responsive, accessibilité, tests et optimisation.
\end{itemize}

\section{Objectifs, enjeux et périmètre}

\subsection{Objectifs fonctionnels}
\begin{enumerate}
  \item \textbf{Module Information} (visiteur) : free tour et recherche d'informations.
  \item \textbf{Module Visualisation} (utilisateur simple) : inscription, connexion, profil, consultation.
  \item \textbf{Module Gestion} (utilisateur complexe) : pilotage des objets, statistiques, maintenance.
  \item \textbf{Module Administration} (administrateur) : gouvernance utilisateurs et supervision globale.
\end{enumerate}

\subsection{Enjeux techniques}
\begin{itemize}
  \item Modéliser un domaine riche avec de nombreux types d'objets connectés.
  \item Assurer une cohérence front/back sur des comportements spécifiques (ex. machine à café, aspirateur, alarmes).
  \item Garantir une démonstration stable en soutenance (données seed réalistes et reproductibles).
\end{itemize}

\section{Technologies et architecture}

\subsection{Stack retenue}
\begin{itemize}
  \item \textbf{Backend} : Java 21, Spring Boot 3.3, Spring Data JPA, Bean Validation.
  \item \textbf{Frontend} : React 19, Vite.
  \item \textbf{BDD} : H2 (mode fichier en développement), préparation MySQL.
  \item \textbf{Déploiement} : Render (API Docker + Frontend statique).
\end{itemize}

\subsection{Architecture logique}
\begin{itemize}
  \item Couche \textbf{modèle} : entités métier (Maison, Pièce, Utilisateur, ObjetConnecte, etc.).
  \item Couche \textbf{repository} : accès aux données via Spring Data.
  \item Couche \textbf{service} : logique métier (points, sessions, scénarios).
  \item Couche \textbf{controller} : API REST par module.
  \item Couche \textbf{frontend} : composants UI modulaires, états locaux, appels API.
\end{itemize}

\subsection{Routes principales}
\begin{itemize}
  \item \texttt{/api/info/*}
  \item \texttt{/api/auth/*}
  \item \texttt{/api/visualisation/*}
  \item \texttt{/api/gestion/*}
  \item \texttt{/api/admin/*}
\end{itemize}

\section{Méthodologie de gestion de projet}

Nous avons appliqué une démarche itérative inspirée de \textbf{SCRUM} :
\begin{itemize}
  \item backlog priorisé (fichiers \texttt{PLAN.md} et \texttt{NEXT\_TASKS.md}),
  \item livraisons incrémentales,
  \item vérifications régulières (builds, tests de parcours, smoke test).
\end{itemize}

\subsection{Répartition des tâches}

\begin{longtable}{|p{3.2cm}|p{5.1cm}|p{6.3cm}|}
\hline
\textbf{Membre} & \textbf{Responsabilité principale} & \textbf{Contributions réalisées} \\
\hline
Khayem Ben Ghorbel & Pilotage technique global & Architecture, intégration front/back, corrections critiques, déploiement, coordination livrables \\
\hline
Rindra & Support conception et implémentation & Contribution aux modules fonctionnels et validation de parcours utilisateur \\
\hline
Ziyad & Support backend/qualité & Aide sur stabilité, validations API et qualité de livraison \\
\hline
Anastasia & UX/présentation & Support structuration interface et lisibilité démo \\
\hline
Valerie & Documentation/soutenance & Support préparation rapport, articulation démonstration et communication \\
\hline
\end{longtable}

\textit{Note : les noms complets de certains membres sont à compléter dans la version finale remise.}

\subsection{Diagramme de Gantt (synthèse)}

\begin{center}
\begin{tabular}{|p{4.6cm}|c|c|c|c|c|c|}
\hline
\textbf{Lot} & \textbf{S1} & \textbf{S2} & \textbf{S3} & \textbf{S4} & \textbf{S5} & \textbf{S6} \\
\hline
Cadrage + architecture & X & X &   &   &   &   \\
\hline
Module Information &   & X & X &   &   &   \\
\hline
Module Visualisation &   &   & X & X &   &   \\
\hline
Module Gestion &   &   &   & X & X &   \\
\hline
Module Administration &   &   &   &   & X & X \\
\hline
Qualité, déploiement, rapport &   &   &   &   & X & X \\
\hline
\end{tabular}
\end{center}

\section{Étapes réalisées (méthodologie de projet)}

\subsection{Phase 1 -- Planification}
\begin{itemize}
  \item Analyse du cahier des charges.
  \item Définition des modules et des rôles utilisateurs.
  \item Choix de la stack technique et de la stratégie de données.
\end{itemize}

\subsection{Phase 2 -- Conception}
\begin{itemize}
  \item Conception du modèle métier UML.
  \item Conception de l'API REST et des DTO.
  \item Structuration des écrans côté frontend.
\end{itemize}

\subsection{Phase 3 -- Développement itératif}
\begin{itemize}
  \item Implémentation incrémentale des modules Information, Visualisation, Gestion, Administration.
  \item Ajout de contrôles spécialisés par type d'objet (ex. machine à café, aspirateur, alarme).
  \item Intégration des scénarios domotiques et du suivi maintenance.
\end{itemize}

\subsection{Phase 4 -- Validation et stabilisation}
\begin{itemize}
  \item Vérification des builds backend et frontend.
  \item Correctifs UX (mise à jour immédiate des états après action).
  \item Construction d'une démo reproductible (seed réaliste + smoke test).
\end{itemize}

\section{Couverture fonctionnelle par module}

\subsection{Module Information}
\begin{itemize}
  \item Consultation des objets et pièces.
  \item Recherche multi-critères (au moins deux filtres).
  \item Entrée utilisateur vers l'inscription.
\end{itemize}

\subsection{Module Visualisation}
\begin{itemize}
  \item Authentification et session utilisateur.
  \item Consultation/édition profil.
  \item Consultation des objets et services avec détails enrichis.
  \item Mécanique de points et progression de niveau.
\end{itemize}

\subsection{Module Gestion}
\begin{itemize}
  \item CRUD objets connectés.
  \item Contrôle d'état, configuration par pièce, paramètres spécifiques.
  \item Statistiques, historique des actions.
  \item Maintenance (détection objets à réparer, marquage réparé).
  \item Exports et supervision opérationnelle.
\end{itemize}

\subsection{Module Administration}
\begin{itemize}
  \item Gestion des utilisateurs et des droits admin.
  \item Traitement des demandes de suppression.
  \item Protection des routes sensibles.
\end{itemize}

\section{UML amélioré}

\subsection{Structure conceptuelle}
Le modèle UML a été enrichi pour couvrir les familles d'objets connectés demandées,
avec héritage et spécialisations métier.

\begin{longtable}{|p{4.2cm}|p{10.4cm}|}
\hline
\textbf{Bloc UML} & \textbf{Rôle} \\
\hline
Maison -- Pièce & Une maison possède plusieurs pièces (Salon, Cuisine, Chambre, etc.) \\
\hline
ObjetConnecte (abstrait) & Racine commune (nom, état, connectivité, batterie, maintenance, pièce) \\
\hline
Branches ObjetConnecte & Ouvrant, Capteur, Appareil, BesoinAnimal \\
\hline
Feuilles métier & Exemples : Porte, Volet, Thermostat, Camera, DetecteurMouvement, LaveLinge, Alarme, Aspirateur, MachineCafe, Nourriture, Eau \\
\hline
Utilisateur (abstrait) & ParentFamille, Enfant, VoisinVisiteur + attributs de niveau/points/admin \\
\hline
HistoriqueAction & Trace les actions utilisateur sur les objets \\
\hline
Scenario + ScenarioAction & Automatisation (manuel/planifié/conditionnel) \\
\hline
DemandeSuppression & Workflow de suppression validé par admin \\
\hline
\end{longtable}

\subsection{Cardinalités clés}
\begin{itemize}
  \item 1 Maison -- N Pièces
  \item 1 Pièce -- N ObjetsConnectés
  \item 1 Utilisateur -- N HistoriqueAction
  \item 1 ObjetConnecte -- N HistoriqueAction
  \item 1 Scenario -- N ScenarioAction
  \item 1 Utilisateur -- N DemandeSuppression
\end{itemize}

\section{Schéma de base de données amélioré}

\subsection{Choix de modélisation relationnelle}
\begin{itemize}
  \item Stratégie d'héritage \texttt{SINGLE\_TABLE} pour \texttt{ObjetConnecte} et \texttt{Utilisateur}.
  \item Clés techniques auto-générées (\texttt{id} bigint).
  \item Clés étrangères explicites pour garantir l'intégrité référentielle.
\end{itemize}

\subsection{Tables principales}

\begin{longtable}{|p{3.1cm}|p{3.6cm}|p{7.8cm}|}
\hline
\textbf{Table} & \textbf{Clé primaire} & \textbf{Contenu principal} \\
\hline
maison & id & Nom de la maison \\
\hline
piece & id & Pièces + FK vers maison \\
\hline
utilisateur & id & Profil, points, niveau, rôle admin, type membre \\
\hline
objet\_connecte & id & Données communes + colonnes spécialisées par type \\
\hline
historique\_action & id & Action, timestamp, FK utilisateur, FK objet \\
\hline
donnee\_capteur & id & Série temporelle de mesures capteurs \\
\hline
scenario & id & Définition scénario, type, trigger, activation \\
\hline
scenario\_action & id & Actions élémentaires d'un scénario, FK scénario, FK objet \\
\hline
demande\_suppression & id & Workflow de demande suppression + statut + auteur \\
\hline
\end{longtable}

\subsection{Améliorations proposées du schéma}
\begin{enumerate}
  \item Indexer les colonnes de recherche fréquentes :
  \begin{itemize}
    \item \texttt{objet\_connecte(type\_objet, piece\_id, etat, marque)}
    \item \texttt{historique\_action(utilisateur\_id, objet\_id, timestamp)}
    \item \texttt{scenario(enabled, type, trigger\_event)}
  \end{itemize}
  \item Contraintes d'unicité utiles :
  \begin{itemize}
    \item \texttt{utilisateur.email} unique,
    \item éventuellement \texttt{objet\_connecte(nom, piece\_id)} unique selon règle métier.
  \end{itemize}
  \item Journaliser les exports/administrations sensibles dans \texttt{historique\_action}.
\end{enumerate}

\section{Tests, performance, responsive, accessibilité}

\subsection{Vérifications techniques}
\begin{itemize}
  \item \texttt{mvn -DskipTests compile} : succès.
  \item \texttt{npm.cmd run build} : succès.
  \item Smoke test fonctionnel : \texttt{scripts/demo-smoke.ps1}.
\end{itemize}

\subsection{Responsive et UX}
\begin{itemize}
  \item Vérifications desktop/tablette/mobile.
  \item Interface orientée feedback (toasts, états lisibles, contrôles par type d'objet).
  \item Navigation modulaire claire (Information / Visualisation / Gestion / Administration).
\end{itemize}

\subsection{Accessibilité (base)}
\begin{itemize}
  \item Structure sémantique globale.
  \item Focus visible et interactions clavier sur les actions principales.
  \item Contrastes et hiérarchie de lecture améliorés.
\end{itemize}

\section{Déploiement et exploitation}

\subsection{Render}
\begin{itemize}
  \item Service API : \texttt{dev-web-ing1-api}
  \item Service Frontend : \texttt{dev-web-ing1-web}
  \item Healthcheck : \texttt{/api/health}
\end{itemize}

\subsection{Important pour la soutenance}
L'URL racine de l'API (\texttt{/}) retourne une 404 Whitelabel.
C'est normal pour une API-only. Le point d'entrée de supervision est
\texttt{/api/health}.

\section{Difficultés rencontrées et solutions}

\begin{enumerate}
  \item \textbf{Évolutions du modèle et compatibilité données} :
  correction de problèmes de valeurs legacy, validations renforcées.

  \item \textbf{Synchronisation front/back en temps réel} :
  mise à jour locale immédiate après \texttt{PUT} pour éviter l'impression de non-réactivité.

  \item \textbf{Fiabilité démo} :
  seed orienté soutenance (objets à réparer garantis, scénarios disponibles, données cohérentes).
\end{enumerate}

\section{Conclusion et perspectives}

Le projet répond au cahier des charges sur les axes techniques, fonctionnels et méthodologiques.
La plateforme est utilisable en démonstration temps réel, structurée pour évoluer et prête à une
version plus industrialisée.

\subsection*{Perspectives}
\begin{itemize}
  \item Passage complet MySQL avec migration contrôlée.
  \item Renforcement des tests automatiques (intégration backend et E2E frontend).
  \item Amélioration de l'observabilité (métriques, logs structurés, audit enrichi).
  \item Renforcement des contrôles de sécurité et des workflows administratifs.
\end{itemize}

\section*{Annexe A -- Checklist de conformité au cahier des charges}

\begin{itemize}
  \item Rapport (<= 15 pages) : présent.
  \item Introduction / contexte / choix techniques : présent.
  \item Répartition des tâches + Gantt : présent.
  \item Étapes réalisées (méthodologie) : présent.
  \item Conclusion + perspectives : présent.
  \item Utilisation Git : présent.
  \item Frameworks : présent.
  \item BDD : présente.
  \item Responsive, accessibilité, tests : documentés.
\end{itemize}

\section*{Annexe B -- Commandes utiles}

\begin{Verbatim}[fontsize=\small]
# backend
mvn spring-boot:run

# frontend
cd frontend
npm.cmd run dev

# build
mvn -DskipTests compile
cd frontend && npm.cmd run build

# smoke demo
powershell.exe -ExecutionPolicy Bypass -File scripts\demo-smoke.ps1
\end{Verbatim}

\end{document}
